\documentclass[english]{letter}
\usepackage[T1]{fontenc}
\usepackage[latin1]{inputenc}
\usepackage{babel}
\usepackage[top=.5in,bottom=1in,left=1in,right=1in,noheadfoot]{geometry}
\usepackage{graphicx}

\address{\sffamily
  \begin{tabular}{rr}
    Prof.~Dr.~Ant�nio Augusto Fr�hlich & \raisebox{-0.9cm}[0pt][0pt]{\resizebox*{!}{1.2cm}{\includegraphics{/home/guto/etc/logos/lisha}}}\\
    Federal University of Santa Catarina &\\
    Software/Hardware Integration Laboratory &\\
    \multicolumn{2}{r}{PO Box 476, 88040-900 Florian�polis, SC, Brazil}\\
    \multicolumn{2}{r}{Phone: +55 48 3721-9516}\\
    \multicolumn{2}{r}{E-Mail: guto@lisha.ufsc.br}\\
    \multicolumn{2}{r}{WWW: http://lisha.ufsc.br/$\sim$guto}\\
    \\
    \multicolumn{2}{r}{\today}
  \end{tabular}
}

\date{}

\signature{
  \begin{center}
   \raisebox{-2em}[0pt][0pt]{\includegraphics[scale=.75]{/home/guto/personal/signature}}\\
   \hrulefill\\
    Ant�nio Augusto Fr�hlich
  \end{center}
}


\begin{document}

\begin{letter}{International Journal of Distributed Sensor Networks\\
	Editor\\
    Prof. Cem Ersoy
}
\opening{Dear Prof. Ersoy,}
\setlength{\parskip}{1.5ex}

A new revised version of the manuscript "A Comprehensive Approach to Power
Management in Embedded Systems", initially submitted to the International
Journal of Distributed Sensor Networks on September 2010 and ressubmited on
January 2011, has just been uploaded through Hindawi manuscript system. Bellow
you can find replies to each individual reviewer.

\textbf{\large Reviewer 4}

\textbf{Remarks}

The author added experimental results of the C-MAC and the AD-ZRP for showing
idea related to WSNs. They show interesting results, but the author just
explains these schemes. How to be integrated these schemes to EPOS? How to be
related with API for saving energy?

\textbf{Reply}

Thank you for your constructive suggestions. Both paragraphs, the one about
C-MAC and the one about AD-ZRP, were rephrased to include explicit references to
the PM strategies discussed along the paper.


\textbf{\large Reviewer 2}

\textbf{Remarks}

The paper was extended by two additional examples showing the suitablility of
the approach for sensor networks. Although the paper is still not FOCUSED on
sensor networks, it shows that the technique may be used within sensor networks.
Unfortunately, the abstract and the conclusions are identical to the original
submission. Both sections should be adjusted according to the additional
examples. Additionally, the reference mentioned by the second reviewer was not
correctly embedded: The approach presented by the authors does not deal with
PCs. As I understood, the approach focuses on embedded soft real-time systems.
The issue that you call "developer knowledge" (i.e. the schedule of several
tasks) is tried to be identified automaically by the approach mentioned in [19].
Hence, you should have a deeper look into that paper and compare it directly to
your approach (they focus only on soft real-time, so there is an obvious
difference but they do not need any additional developers knowledge).

\textbf{Reply}

Thank you for your constructive suggestions. The decision not to narrow the
paper's scope to focus only on WSN was explained in the previous reply,
however, I must apology for not having adjusted abstract and conclusion
accordingly. This new version addresses these two points explicitly.

In respect to the approach presented in [19], I believe that, just like the
proposals in Grace-OS, ECOS, and Niu's (m,k)-scheduler, it can be deployed
together with the mechanisms presented here straightforwardly. In the revised
version, I stress I/O as the critical component that is missing in many other
related works. DVS-based approaches assume regular (or somehow predictable)
execution patterns and assume also that CPU time can be directly correlated to
energy consumption. From WSN, for example, we know that channel contention may
lead to energetically expensive cycles that cannot be accounted from their
execution time. Application's exception handling (like the trigger thread in the
paper) is also usually hard to account from CPU usage. This is explained across
section IV: initially in the second paragraph, then on the third topic of the
subsequent bullet list and finally through related works. Explicitly discussing
the autocorrelation clustering DVS technique proposed in [19] would fit better
in a paper about DVS-based scheduling. Here we sustain that both aspects of
power management are needed and must be deployed together (e.g., in the top
paragraph, page 10, second column).


\textbf{\large Reviewer 1}

I apology for mistyping Petri Nets. It has been fixed.

Sincerely,

\closing{~}

\end{letter}

\end{document}





















































